\section{Conclusion}
This paper asked whether Sentinel-2 imagery stored on the HEALPix level next
finer than its native resolution, which is necessary to avoid loss of
spatial detail but nominally oversamples the data by a factor of
$\approx2.58$ in cell area, costs more to store once compression and the
storage layout of the native product are both taken into account. Using
only established, independently published compressors (ZFP, SZ3, JPEG~2000,
JPEG~XL and a DWT+Zstandard baseline), no HEALPix-specific coding, and an
equal-area protocol that crops the native UTM raster to a square of the same
nominal ground area as each HEALPix tile, the answer on 40 regions spanning
four land-cover classes and both HEALPix cell-shape regimes is that native
UTM is intrinsically the more compressible representation throughout, and
that whether the carrier becomes operationally competitive depends on the
tolerated distortion and on the overhead of the native layout.

Before compression, the carrier tile is by construction $2.584\times$ the
equal-area UTM crop. After compression, the intrinsic HEALPix/UTM gap at
$\gamma=1$ decreases monotonically with looser targets, from $+106\%$ to
$+6\%$, which corresponds to a critical overhead $\gamma_\star$ falling from
$2.06$ to $1.06$. This decrease is a property of the sampling geometry
rather than of scene content: an analytic field with the spectrum and
instrument width fitted to the imagery, sampled by both lattices with no
support mismatch and no content difference, reproduces the gap to within
$8$--$19\%$ at every target. At the illustrative overhead
$\gamma_{\mathrm{tile}}=1.30$, native UTM remains smaller at the three
tightest RMSE targets and the two looser budgets (median ratios
$1.23$--$1.59$ and $1.33$--$4.02$), HEALPix is smaller at the loosest RMSE
target ($0.81$) and the tightest budget ($0.83$), and seven of the eight
paired tests survive the Holm correction ($p\le8\times10^{-3}$, $n=40$). The
compressed size of HEALPix falls below the operational comparison only where
the real overhead of the native product exceeds $\gamma_\star$, which is a
statement about product layout rather than about UTM as a coordinate
system.

The crossover is a result for one transfer operator and does not survive
the others tested. The carrier was populated with Clough--Tocher
interpolation, chosen before any storage figure was measured
(Section~\ref{sec:resampler}); rebuilding the same tiles with a PSF-aware or
an inverse-distance transfer costs $1.2$ to $2.2$ times the bytes at equal
accuracy and leaves HEALPix ahead at no target
(Section~\ref{sec:transfer_cost}). Only a few percent of this cost is the
price of deconvolution; the remainder is the price of leaving a $C^1$
curvature-minimising interpolant, whose output a lossy coder can discard far
more of at equal error, and the observations are consistent with energy
compaction rather than total information content as the mechanism. The
observed crossover is therefore a result for a carrier populated by a
transfer in the smoothness class represented here by Clough--Tocher, and
should not be read as a transfer-independent property of HEALPix. On the
codec side, JPEG~2000 and JPEG~XL lead on both grids alike, which we
interpret as a property of these codecs on this imagery rather than of the
quadtree.

The verdict also depends on which reconstruction reference is accepted, and
both are reported. Against the same Sentinel-2 observation, the uncompressed
carrier already differs from the native product by $4.0\times10^{-3}$ in
median under the nearest-cell read, about $10\%$ of the standard deviation
of a scene, and no swept configuration reaches four of the five RMSE targets
against it; the $2\%$-budget ratio changes from $0.83$ to $1.44$ in favour
of native UTM. This floor belongs to the transfer-and-read pair rather than
to the grid: a read that interpolates within the carrier lattice removes two
thirds of its variance on the same tiles (Table~\ref{tab:reader_floor}), and
when the sweep is re-scored through that read the lead of native UTM narrows
to $1.12$ without reversing, the loosest RMSE target reaches parity ($1.02$,
$n=40$), and the residual floor of $1.9\times10^{-3}$ is due to the
transfer. Read against the latent scene instead, the reference that applies
when the carrier is used as if it were the scene, HEALPix becomes the
smaller representation while the compression error is still $0.13$ to
$0.50$ of the error the product already carries, in three of the four scene
classes (Appendix~\ref{app:error_budget}). The practical conclusion is
conditional and should be stated as such: a HEALPix carrier is competitive
at aggressive compression when the HEALPix cube is itself the reference
product, and is not when the native product remains the ground truth and
the transfer error is charged to it.

In summary, an oversampled HEALPix carrier does not intrinsically compress
Sentinel-2 imagery more efficiently than its native UTM representation, but
its oversampling penalty decreases under lossy compression, and whether it
becomes operationally competitive is determined by the tolerated
distortion, the relative storage overhead of the two layouts, the transfer
used to build the carrier and the reference used to evaluate it. The
critical-overhead formulation expresses this trade-off without privileging
any storage architecture, and it is the form in which we consider the
result should be used. The evidence is from Sentinel-2 B04 at \SI{10}{m},
one acquisition per region, and should not be transferred to multispectral
products, temporal cubes, model output or flux-conservative fields without
being re-measured there. The most important next step is to decouple the
three parts of the chain that this paper has shown to matter, namely the
transfer, the carrier and the reading operator, and to evaluate the
complete chain with the operational PSF-aware transfer and a reader matched
to it on the same sites, which is what the pipeline shared with the
companion study~\cite{delouis2026resample} was built for.
